\chapter{Predicate Logic: Syntax}\label{ch:predsyntax}

\section{Why propositional logic is not enough}

Consider the argument:
\begin{quote}
Every student in CS 245 is clever. Alice is a student in CS 245. Therefore Alice is clever.
\end{quote}
In propositional logic these are three unrelated atoms $p,q,r$, and $\{p,q\}\not\models r$.
The validity depends on the internal structure of the sentences: individuals (Alice),
properties (``is clever''), and quantifiers (``every''). First-order predicate logic adds
exactly these:
\[
  \forall x\,(S(x)\to C(x)),\ S(a)\ \models\ C(a).
\]

\section{The language}

A first-order language $\mathcal{L}$ consists of logical symbols, common to all
languages, and non-logical symbols, which depend on the application.

\paragraph{Logical symbols.}
\begin{itemize}
  \item connectives $\neg,\land,\lor,\to,\leftrightarrow$ and parentheses, comma;
  \item quantifiers $\forall$ (for all) and $\exists$ (there exists);
  \item an infinite supply of \emph{variables} $x,y,z,u,v,x_1,\dots$;
  \item (optionally) the equality symbol $=$.
\end{itemize}

\paragraph{Non-logical symbols (the signature).}
\begin{itemize}
  \item \emph{constant symbols} $a,b,c,\dots$ (names for individuals);
  \item \emph{function symbols} $f,g,h,\dots$, each with an arity $n\ge1$;
  \item \emph{predicate (relation) symbols} $P,Q,R,\dots$, each with an arity $n\ge 0$.
        A $0$-ary predicate is a propositional atom.
\end{itemize}
(Constants can be viewed as $0$-ary function symbols.)

\begin{example}[Language of arithmetic]
Constant $0$; unary function $s$ (successor); binary functions $+$ and $\cdot$ (written
infix); binary predicate $<$; and equality.
\end{example}

\section{Terms}

Terms denote individuals.
\begin{definition}[Term]
The set of terms is the smallest set such that
\begin{enumerate}
  \item every variable and every constant symbol is a term;
  \item if $f$ is an $n$-ary function symbol and $t_1,\dots,t_n$ are terms, then
        $f(t_1,\dots,t_n)$ is a term.
\end{enumerate}
A term with no variables is a \emph{closed} (ground) term.
\end{definition}

\begin{example}
In arithmetic, $0$, $x$, $s(s(0))$, $(x+s(0))\cdot y$ are terms. $s(0)<x$ is \emph{not} a
term (it is a formula), and neither is $s$ alone.
\end{example}

\section{Formulas}

\begin{definition}[Atomic formula]
If $P$ is an $n$-ary predicate symbol and $t_1,\dots,t_n$ are terms, then
$P(t_1,\dots,t_n)$ is an atomic formula. If $=$ is present, $t_1 = t_2$ is also an atomic
formula.
\end{definition}

\begin{definition}[Formula]
The set of formulas of $\mathcal{L}$ is the smallest set such that
\begin{enumerate}
  \item every atomic formula is a formula;
  \item if $\varphi,\psi$ are formulas then so are $(\neg\varphi)$ and $(\varphi\star\psi)$
        for $\star\in\{\land,\lor,\to,\leftrightarrow\}$;
  \item if $\varphi$ is a formula and $x$ a variable, then $(\forall x\,\varphi)$ and
        $(\exists x\,\varphi)$ are formulas.
\end{enumerate}
\end{definition}

Quantifiers bind as tightly as $\neg$: $\forall x\,P(x)\to Q(x)$ means
$(\forall x\,P(x))\to Q(x)$. Use parentheses when the scope should extend further.

As in the propositional case, unique readability holds and we can draw parse trees:
quantifier nodes $\forall x$ and $\exists x$ have one child, predicate nodes have their
argument terms as children.

\begin{important}[Terms vs.\ formulas]
Terms (built from variables, constants, functions) denote \emph{objects}; formulas (built
from predicates, connectives, quantifiers) denote \emph{truth values}. You cannot apply a
function to a formula or a connective to a term: $f(P(x))$ and $\neg x$ are not
well-formed.
\end{important}

\section{Translating English}

\begin{itemize}
  \item ``All $A$s are $B$s'': $\forall x\,(A(x)\to B(x))$ \quad(\emph{not} $\land$).
  \item ``Some $A$ is a $B$'': $\exists x\,(A(x)\land B(x))$ \quad(\emph{not} $\to$).
  \item ``No $A$ is a $B$'': $\forall x\,(A(x)\to\neg B(x))$, or
        $\neg\exists x\,(A(x)\land B(x))$.
  \item ``Not all $A$s are $B$s'': $\exists x\,(A(x)\land\neg B(x))$.
  \item ``There is exactly one $x$ with $P$'':
        $\exists x\,(P(x)\land\forall y\,(P(y)\to y=x))$.
  \item ``There are at least two'': $\exists x\,\exists y\,(\neg(x=y)\land P(x)\land P(y))$.
\end{itemize}

\begin{important}[Pitfall: $\land$ vs.\ $\to$ with quantifiers]
$\forall x\,(S(x)\land C(x))$ says \emph{everything} is a clever student.
$\exists x\,(S(x)\to C(x))$ is true as soon as there is any non-student, so it says
almost nothing. Pair $\forall$ with $\to$ and $\exists$ with $\land$.
\end{important}

\begin{example}[Order of quantifiers]
With $L(x,y)$ meaning ``$x$ loves $y$'':
\begin{itemize}
  \item $\forall x\,\exists y\,L(x,y)$: everybody loves somebody (possibly different
        people);
  \item $\exists y\,\forall x\,L(x,y)$: there is one person whom everybody loves.
\end{itemize}
The second implies the first but not conversely.
\end{example}

\begin{example}
``Every natural number has a larger prime'':
$\forall x\,\exists y\,(x<y\land \mathit{Prime}(y))$, where $\mathit{Prime}(y)$ could
itself be defined as
$s(0)<y\land\forall u\,\forall v\,(u\cdot v=y\to(u=s(0)\lor v=s(0)))$.
\end{example}

\section{Free and bound variables}

\begin{definition}[Scope, free and bound occurrences]
In $\forall x\,\varphi$ or $\exists x\,\varphi$, the formula $\varphi$ is the \emph{scope}
of the quantifier. An occurrence of $x$ is \emph{bound} if it lies within the scope of a
quantifier $\forall x$ or $\exists x$ (or is the $x$ right after the quantifier);
otherwise it is \emph{free}.
\end{definition}

The set of free variables is defined recursively:
\begin{align*}
  FV(P(t_1,\dots,t_n)) &= \mathrm{vars}(t_1)\cup\dots\cup\mathrm{vars}(t_n),\\
  FV(\neg\varphi) &= FV(\varphi),\qquad FV(\varphi\star\psi)=FV(\varphi)\cup FV(\psi),\\
  FV(\forall x\,\varphi) &= FV(\exists x\,\varphi) = FV(\varphi)\setminus\{x\}.
\end{align*}

\begin{definition}[Sentence]
A formula with no free variables is a \emph{sentence} (closed formula). Only sentences have
a truth value once an interpretation is fixed; a formula with free variables also needs
values for those variables.
\end{definition}

\begin{example}
In $\forall x\,(P(x,y)\to\exists y\,Q(y,x))\land R(x)$:
\begin{itemize}
  \item the $x$ in $P(x,y)$ and in $Q(y,x)$ are bound by $\forall x$;
  \item the $y$ in $P(x,y)$ is \emph{free}; the $y$ in $Q(y,x)$ is bound by $\exists y$;
  \item the $x$ in $R(x)$ is free (outside the scope of $\forall x$).
\end{itemize}
So $FV = \{x,y\}$. The same variable can occur both free and bound in one formula.
\end{example}

Bound variables are ``dummy'' names: $\forall x\,P(x)$ and $\forall y\,P(y)$ say the same
thing (they are \emph{alpha-equivalent}), just as $\sum_i a_i = \sum_j a_j$.

\section{Substitution}

\begin{definition}[Substitution]
$\varphi[t/x]$ denotes the result of replacing every \emph{free} occurrence of $x$ in
$\varphi$ by the term $t$. Formally, on terms $x[t/x]=t$, $y[t/x]=y$ for $y\ne x$, and
$f(t_1,\dots)[t/x]=f(t_1[t/x],\dots)$; on formulas it passes through connectives, and
\[
  (\forall y\,\varphi)[t/x] =
  \begin{cases}
   \forall y\,\varphi & \text{if } y=x \ (\text{$x$ is not free}),\\
   \forall y\,(\varphi[t/x]) & \text{if } y\ne x,
  \end{cases}
\]
and similarly for $\exists$.
\end{definition}

\begin{definition}[Free for]
A term $t$ is \emph{free for} $x$ in $\varphi$ if no free occurrence of $x$ in $\varphi$ lies
within the scope of a quantifier $\forall y$ or $\exists y$ for some variable $y$ occurring
in $t$. In words: substituting $t$ does not cause any variable of $t$ to be
\emph{captured}.
\end{definition}

\begin{example}[Variable capture]
Let $\varphi = \exists y\,(x<y)$ (``something is bigger than $x$''; true in $\N$ for every
$x$). Then
\begin{itemize}
  \item $\varphi[s(z)/x]=\exists y\,(s(z)<y)$: fine, $s(z)$ is free for $x$.
  \item $\varphi[y/x]=\exists y\,(y<y)$: wrong! The $y$ got captured, and the meaning
        changed to ``something is bigger than itself'' (false). $y$ is not free for $x$.
\end{itemize}
The fix is to rename the bound variable first: $\exists y\,(x<y)\equiv\exists w\,(x<w)$, so
$\varphi[y/x]$ should be taken as $\exists w\,(y<w)$.
\end{example}

\begin{important}[Substitution rules of thumb]
\begin{itemize}
  \item Only \emph{free} occurrences are replaced.
  \item Before substituting $t$ for $x$, rename any bound variables of $\varphi$ that occur
        in $t$ and whose quantifiers have free $x$ in their scope.
  \item Every rule in Chapter~\ref{ch:preddeduction} that substitutes ($\forall e$,
        $\exists i$) requires $t$ to be free for $x$.
\end{itemize}
\end{important}

\section*{Summary}
\begin{itemize}
  \item A first-order language has constants, function symbols and predicate symbols,
        plus variables, connectives, quantifiers and possibly $=$.
  \item Terms denote objects; atomic formulas apply predicates to terms; formulas close
        under connectives and quantifiers.
  \item Translate ``all $A$ are $B$'' with $\forall$ and $\to$; ``some $A$ is $B$'' with
        $\exists$ and $\land$; quantifier order matters.
  \item Free vs.\ bound occurrences; sentences have no free variables.
  \item $\varphi[t/x]$ replaces free $x$ by $t$; it is only safe when $t$ is free for $x$
        (no capture); rename bound variables otherwise.
\end{itemize}
